\section{Implications and Predictions}

If the picture holds, several things we are used to treating as separate turn out to be one, and a few sharp consequences become testable. The implications are worth stating plainly, because they are the reason the model is worth the trouble, and the predictions are worth stating plainly because they are how it could be shown wrong.

The deepest implication is a set of collapses, where three or four ideas that physics keeps apart are revealed as one. Time, the arrow of time, and the expansion of space are not three facts but one, the growing of the crystal, felt as duration, pointing as a direction, and measured as a scale factor. Grand unification is not a hoped-for symmetry to be discovered but the symmetry of the twenty-four directions, read at the cusp, so the Standard Model is geometry rather than a separate input. The world is exactly discrete and finite at every beat, hence exactly computable, with the continuum a coarse-grained appearance rather than a fundamental fact. And consciousness is not a late and puzzling product of complex matter but the substance of the substrate itself, with matter the outer face of a feeling that was there from the first beat. Each of these is a payoff of the single choice of substrate, and together they are why a spare model is offered for so much.

A theory that explained everything and forbade nothing would not be worth having, so the model names its own executioners. It forces exactly three generations of matter, so a clean fourth generation would falsify it. It rests on the grand-unified relations, so a failure of the weak-angle running or of bottom-tau unification at the unification scale would break it. It predicts low-energy supersymmetry as the spectrum that makes the running land correctly, so the permanent absence of any such spectrum would count against it. And it needs the Skyrme stabilizer positive for matter to form, which is measured but could in principle have come out otherwise.

\begin{result}[The cleanest falsifier]
A discrete substrate cannot make Lorentz invariance exact to all orders. It must leave a faint cubic anisotropy, a direction-dependence in the speed of very high-energy signals that scales as the square of the energy in Planck units. The effect is far below current reach and consistent with the tightest gamma-ray-burst bounds, but it is not zero. If Lorentz invariance is ever shown exact to all orders, with no anisotropy at any scale, the substrate is wrong. This is the sharpest experimental handle the model offers.
\end{result}

These are honest stakes. The model makes ratios it gets right, names the absolute numbers it does not yet get, and points to a specific high-energy effect that would settle the matter against it. That is the most a feasibility-stage theory can do, and the paper does it rather than hiding behind the parts that already work.
